\section{Discussion}
\paragraph{An evidence requirement, not a universal winning architecture.}
The benchmark exposes distinctions that an individual-event leaderboard
cannot resolve. A shared factor can improve collective dispersion without
changing a single individual probability. Typed incidence can then add a
smaller, projection-specific gain. Correcting a common bias or forecasting
counts directly can change the ranking again. These are empirical reasons to
include global, direct-count and calibration controls before attributing
collective accuracy to a temporal graph architecture.

\paragraph{Which output is needed?}
If the application asks only for a few fixed group counts, direct count models
are strong and economical candidates. Entity-level factor scenarios supply a
different capability: the same realization can be projected onto overlapping
jobs and machines or an alternative grouping without fitting another count
model. A complete handler hierarchy can also be reconciled coherently. Neither
capability guarantees better prediction. In particular, separate good group
marginals need not give the right extreme across groups, and a globally
coherent scenario can still have biased means or incorrect tails. The joint
profile and peak-deviation scores therefore complement count CRPS.

\paragraph{Calibration and dependence interact.}
Replacing biased individual probabilities changes both expected group counts
and the quantile needed by a reserve proxy. A coupling that improved CRPS
before correction need not have the same decision effect afterward. The
same-version requirement matters because a sequence of chronological OOF
predictors is not one stationary calibration object. Nevertheless, delayed
updates can lag a regime change, and the fresh replication does not justify
claiming that rolling correction always beats a fixed earlier map. A factor
estimated from forecast residuals may absorb misspecification; its value is
predictive rather than evidence of an identified common cause.

\paragraph{Limits of generalization.}
The evidence covers one compute cluster and two workflow releases, with only
seven final Borg days and 28 fresh workflow origins. Four or two bootstrap
blocks in the latter are weak support for long-run uncertainty. Repeated cases,
millions of tasks and Monte Carlo seeds do not multiply independent temporal
replications. Retrospective follow-ups are separated from their motivating
outcomes, and the single fresh confirmation should not be interpreted as
broad external validation. No live operator study, policy deployment, measured
cost saving or resource-capacity ground truth is available.

The benchmark conditions on the currently active cohort and excludes future
arrivals, subsequent exits/returns and membership evolution. Last handler is
an observed association, not exclusive ownership. Positive factor families
cannot represent general negative dependence. A full crossed job--machine
probabilistic reconciliation baseline, adaptations of TACTiS or graph
point-process generators, and complete business-process simulation remain
unimplemented comparisons. The paper makes no superiority claim over them.
The three individual architectures and fixed training budgets are useful
controls, not an exhaustive model search.

\paragraph{Implications for scientific dataset generation.}
A reusable scenario source should declare the conditioning population,
observable history and outcome contract before generating records. It should
also test both individual and collective statistics under meaningful group
projections. Our adapters, score suite and provenance checks provide a bounded
foundation for that purpose. They are not a full synthetic-data factory and
have not been validated by training downstream models on synthetic data.
Extending them requires arrivals, repeated transitions and separate utility
and privacy studies. Original public-source license and attribution terms
remain applicable.
